\section{Experiment Setup}
\label{sec:experiment-setup}

% 介绍模型使用情况：一个用于微调，一个用于垂域微调对比，另有第三方大模型。
\subsection{Models}

We use Qwen3-4B~\citep{yang-etal-2025-qwen3} as the primary backbone and fine-tune 
it separately on each scientific-writing evaluation task. For comparison, we 
evaluate SciRM~\citep{sahinuc-etal-2026-reward}, a reward model developed 
specifically for scientific-writing assessment, as well as several third-party 
large language models under the same evaluation protocol.
%todo：还需要参考 reward modeling 论文介绍用的三方大模型去修改

% 引入原始数据来源和情况介绍
\subsection{Dataset}

Following \citet{sahinuc-etal-2026-reward}, we construct the evaluation suite from 
two complementary sources. The related-work tasks are based on the fine-grained 
evaluation framework introduced by \citet{sahinuc-etal-2025-related-work} and 
extended by \citet{sahinuc-etal-2026-reward} to cover additional machine-learning 
venues. The review-utility tasks are drawn from 
RevUtil~\citep{sadallah-etal-2025-good} and focus on the quality of peer-review 
comments.

\subsubsection*{Related-work data}

The related-work data evaluate literature synthesis along three dimensions:
coherence, positioning type, and positioning check. Each task is formulated
as a binary classification problem, with score 0 indicating that the
specified criterion is not satisfied and score 1 indicating that it is
satisfied.

\begin{enumerate}
    \renewcommand{\labelenumi}{(\arabic{enumi})}

    \item \textit{coherence}~\citep{sahinuc-etal-2025-related-work}:
    Determines whether a citation sentence accurately reflects and remains
    consistent with the findings of the cited paper.

    \item \textit{positioning type}~\citep{sahinuc-etal-2025-related-work}:
    Determines whether the related-work section presents the paper's
    contribution or positioning in the required location and form.

    \item \textit{positioning check}~\citep{sahinuc-etal-2025-related-work}:
    Determines whether the related-work section compares the target paper
    with prior studies and establishes its position in the literature.

\end{enumerate}

\subsubsection*{Review-utility data}

The review-utility data evaluate peer-review comments along four dimensions:
actionability, grounding specificity, helpfulness, and verifiability. These
tasks use five-point ordinal scales, with scores from 1 to 5 and higher scores
indicating more useful, specific, actionable, and better-supported feedback.

\begin{enumerate}
    \renewcommand{\labelenumi}{(\arabic{enumi})}

    \item \textit{actionability}~\citep{sadallah-etal-2025-good}: Measures
    how explicitly and concretely a review comment specifies the changes that
    authors should make.

    \item \textit{grounding specificity}~\citep{sadallah-etal-2025-good}:
    Measures whether a review comment identifies the relevant part of the
    paper and specifies the issue that should be addressed.

    \item \textit{helpfulness}~\citep{sadallah-etal-2025-good}: Measures
    whether a review comment provides clear, constructive, and actionable
    guidance for improving the paper.

    \item \textit{verifiability}~\citep{sadallah-etal-2025-good}: Measures
    whether the claims in a review comment are supported by sufficient
    reasoning, evidence, or external references.

\end{enumerate}

The processing of the source data and the construction of the final training
and test sets are detailed in Appendix~\ref{app:dataset-process}.
Representative examples from the resulting datasets are provided in
Appendix~\ref{app:dataset-examples}.
